% =============================================================================
% Theory — Forward Model Summary
% =============================================================================
\section{Theory}\label{sec:theory}

\subsection{Shell geometry and material model}

We model the organ as an oblate spheroidal shell with semi-major axis~$a$,
semi-minor axis~$c < a$, uniform wall thickness~$h$, and constitutive
properties described by Young's modulus~$E$, Poisson's ratio~$\nu$, wall
density~$\rho_w$, and a Kelvin--Voigt loss tangent~$\eta$.  The cavity is
filled with an incompressible fluid of density~$\rho_f$ and bulk
modulus~$K_f$.  A static intra-abdominal pressure~$P_\text{iap}$ pre-stresses
the membrane.

For the canonical human abdomen, $a = \SI{0.18}{m}$, $c = \SI{0.12}{m}$,
$h = \SI{0.010}{m}$, $E = \SI{0.1}{MPa}$, $\nu = 0.45$,
$\rho_w = \SI{1100}{kg/m^3}$, $\rho_f = \SI{1020}{kg/m^3}$,
$K_f = \SI{2.2}{GPa}$, $P_\text{iap} = \SI{1000}{Pa}$, and $\eta = 0.25$.

\subsection{Equivalent-sphere model}

The simplest forward model replaces the oblate spheroid with an equivalent
sphere of radius $R_\text{eq} = (a^2 c)^{1/3}$.  The flexural frequency of
mode~$n$ ($n \geq 2$) is then
\begin{equation}\label{eq:sphere-freq}
  \omega_n^2
  = \frac{K_\text{bend}(n) + K_\text{memb}(n) + K_\text{pre}(n)}
         {m_\text{wall}(n) + m_\text{fluid}(n)},
\end{equation}
where all stiffness and mass terms depend on $R_\text{eq}$, $h$, $E$, $\nu$,
$\rho_w$, $\rho_f$, and $P_\text{iap}$ through well-known shell-theory
expressions.

Crucially, because every term in \eqref{eq:sphere-freq} involves only the
\emph{single} length scale~$R_\text{eq}$, the ratio $f_n / f_m$ for any two
modes $n, m$ is independent of $(a, c)$ --- it depends only on shell-theory
mode-shape factors.  This means that geometry and stiffness are
\emph{confounded}: many different $(a, c, E)$ triples produce the same
spectrum.

\subsection{Rayleigh--Ritz oblate model}\label{sec:ritz}

The Ritz model retains the oblate geometry.  For each mode number~$n$, we
adopt a two-degree-of-freedom trial displacement:
\begin{equation}
  w(\eta) = \beta \, P_n(\eta), \qquad
  u(\theta) = \alpha \, \sin\theta \, P_n'(\cos\theta),
\end{equation}
where $P_n$ is the Legendre polynomial and $(\alpha, \beta)$ are generalised
coordinates.  The strain energy and kinetic energy are assembled by
Gauss--Legendre quadrature over the actual oblate surface, yielding a $2
\times 2$ generalised eigenvalue problem
\begin{equation}
  \mathbf{K} \mathbf{q} = \omega^2 \mathbf{M} \mathbf{q},
\end{equation}
whose smallest positive eigenvalue gives $\omega_n$.

Because the oblate surface has position-dependent curvature, different modes
``sample'' the geometry differently: the meridional and circumferential
curvatures at the pole differ from those at the equator.  As a result, the
frequency ratios $f_n / f_m$ \emph{do} depend on the aspect ratio~$a/c$,
breaking the confounding seen in the sphere model.

\subsection{The inverse problem}

Let $\boldsymbol{\theta} = (a, c, E)^\top$ denote the parameter vector and
$\mathbf{f}(\boldsymbol{\theta}) = (f_2, f_3, \dots, f_{N+1})^\top$ the
vector of flexural frequencies.  The inverse problem is: given
observed~$\mathbf{f}^\text{obs}$, find~$\boldsymbol{\hat\theta}$ such that
$\mathbf{f}(\boldsymbol{\hat\theta}) \approx \mathbf{f}^\text{obs}$.

This is well-posed if the Jacobian
\begin{equation}\label{eq:jacobian}
  J_{ij}
  = \frac{\partial f_i}{\partial \theta_j}
\end{equation}
has full column rank.  We quantify identifiability via the condition
number~$\kappa(\mathbf{J})$ and the Fisher information matrix.
